% P25 | 4.9 Hedging | Result: hedging with the wrong volatility
% guide: dev_notes/paragraph_guides/P25_result_hedging_with_the_wrong_volatility.txt
If the stock has realised volatility $\sigma_{real}$ and the hedge uses $\sigma_h$, the discounted profit of the seller is, to leading order,
\begin{equation}\label{eq:gammapnl}
 \varepsilon\approx\int_0^Te^{-rt}\,\tfrac12\,\Gamma_t S_t^2\,(\sigma_h^2-\sigma_{real}^2)\,dt,
\end{equation}
so the seller loses when the realised variance exceeds the variance used to price and hedge.
Table~\ref{tbl:hedgewrong} compares the simulated mean profit with the integral in Equation~\eqref{eq:gammapnl} evaluated along the same paths.
At a realised volatility of $55\%$ the simulation gives \$\WrongVolLoss\ and the formula gives \$\WrongVolTheory.
At realised volatilities of $25\%$ and $30\%$ the sign reverses and the seller gains by comparable amounts, so the profit is close to linear in the variance gap.
The size of the error is large compared with the discretisation error of Table~\ref{tbl:hedgefreq}: choosing the wrong volatility costs far more than trading weekly instead of daily.
Section~\ref{sec:misspec} repeats this experiment for the American put, where a wrong volatility also corrupts the exercise rule.
